% Peer Review Should Be Scaled and Continuous for Every Paper
% Copyright (c) 2026 Rick Mathews. All rights reserved.
% Generated from _essays/peer-review-scaled-and-continuous.md by tools/essay-pdf/md2tex.py; edit the Markdown, not this file.
% Build: latexmk -xelatex <this file>
\documentclass[12pt]{article}
\usepackage[T1]{fontenc}
\usepackage{lmodern}
\usepackage[final]{microtype}
\usepackage[letterpaper,top=0.7in,bottom=0.75in,left=1in,right=1in]{geometry}
\usepackage{ragged2e}
\usepackage{xcolor}
\usepackage{enumitem}
\usepackage[hidelinks]{hyperref}
\usepackage{hyperxmp}
\hypersetup{
  pdftitle={Peer Review Should Be Scaled and Continuous for Every Paper},
  pdfauthor={Rick Mathews, PhD},
  pdfsubject={Essay},
  pdfkeywords={peer review, preprints, research integrity, scientific publishing, essay},
  pdfcopyright={Copyright (C) 2026 Rick Mathews. All rights reserved.},
  pdfcontacturl={https://cortriatriatum.github.io/me/},
  pdflang={en-US},
  pdfcreator={Rick Mathews},
  pdfproducer={LaTeX}
}
\definecolor{rule}{gray}{0.6}
\definecolor{quote}{gray}{0.3}
\pagestyle{empty}
% Ragged right: lines run the full measure without stretched word spacing.
\setlength{\RaggedRightParindent}{1.5em}
\setlength{\RaggedRightRightskip}{0pt plus 2em}
\setlength{\emergencystretch}{1.5em}
\tolerance=1000
\hyphenpenalty=1000 % no words split across lines
\widowpenalty=10000 % no single lines stranded at the top or bottom of a page
\clubpenalty=10000
\begin{document}
\special{pdf:docinfo << /Author (Rick Mathews, PhD) >>}% hyperxmp drops the Author field
\RaggedRight
\begin{center}
  {\fontsize{21}{25.2}\selectfont\bfseries Peer Review Should Be Scaled and Continuous for Every Paper\par}
  \vspace{9pt}
  {\large Rick Mathews, PhD\par}
  \vspace{3pt}
  {\small November 17\textsuperscript{th}, 2026\par}
\end{center}
\vspace{2pt}
{\color{rule}\noindent\rule{\textwidth}{0.4pt}}\par
\vspace{-4pt}

\fontsize{13.0}{18.85}\selectfont
\setlength{\parskip}{6.5pt}
Formal peer review, as practiced today, is significantly flawed. Perhaps chief among its problems is that we treat the judgment of two or three people, made at a single point in time, as a permanent label of quality. At a minimum, we must scale the number of peer reviewers and the length of time a paper undergoes peer review.

The scale of plagiarism, falsification, fabrication, and irreproducibility is evidence for how little peer review verifies. About 2 percent of scientists admit to fabricating or falsifying data\textsuperscript{\hyperlink{ref-fanelli2009}{1}}, and two-thirds of retracted papers were pulled for misconduct, not honest error\textsuperscript{\hyperlink{ref-fang2012}{2}}. When researchers repeated 100 published psychology studies, only about a third held up\textsuperscript{\hyperlink{ref-osc2015}{3}}. When the \textit{BMJ} planted nine major errors in a paper, reviewers caught fewer than three on average\textsuperscript{\hyperlink{ref-schroter2008}{4}}. Peer review was never designed to catch fraud; reviewers rarely see the raw data. That is exactly the problem: we give its stamp a weight it cannot carry.

The economics of peer review are strange too, to say the least. Scientists are paid primarily by public funders such as the NIH. That money pays for the salaries, machines, services, reagents, and employees needed to do research, and the output is a scientific paper. The paper must be published to count as a complete output, so scientists often pay journals, sometimes thousands of dollars, to publish their own work: authors paid the five largest publishers about a billion dollars in open-access fees from 2015 to 2018 alone\textsuperscript{\hyperlink{ref-butler2023}{5}}. Journals then charge universities and libraries subscription or pay-per-view fees to read it. \href{https://www.theguardian.com/science/2017/jun/27/profitable-business-scientific-publishing-bad-for-science}{\textit{The Guardian}} has described how publishers take unpaid labor from scientists, who write and review the papers, and sell the results back to those same scientists' institutions\textsuperscript{\hyperlink{ref-buranyi2017}{6}}. Five publishers now account for more than half of all papers\textsuperscript{\hyperlink{ref-lariviere2015}{7}}, and researchers donated more than 100 million hours to peer review in 2020 alone\textsuperscript{\hyperlink{ref-aczel2021}{8}}. Scientists should be paid for their own work, not pay again to read it, yet academia still ties careers to h-indices, publication counts, and impact factors, none of which truly measure research impact\textsuperscript{\hyperlink{ref-koltun2021}{9}}. Meanwhile, editors struggle every day to find two or three reviewers, who may not be the best suited to the task, and the label those reviewers confer is often unjustified.

For peer review, at least, a scaled and continuous alternative is workable today. Papers posted to preprint servers such as \href{https://www.biorxiv.org/}{bioRxiv} could invite comment and record consensus from a verifiably informed group of peers. Each paper would become a living, breathing document, reviewed continuously by a scaled and informed crowd rather than a few overworked strangers. There is early evidence this works. The chemistry journal \textit{Synlett} gave a crowd of about 100 reviewers 72 hours to comment on submitted papers\textsuperscript{\hyperlink{ref-ritter2017}{10}}, and crowd review is now its default method\textsuperscript{\hyperlink{ref-vangemmeren2020}{11}}. Readers already catch what reviewers miss: screening published papers, one team found problematic images in 3.8 percent of more than 20,000 biomedical papers\textsuperscript{\hyperlink{ref-bik2016}{12}}.

Pieces of this already exist. \href{https://pubpeer.com/}{PubPeer} hosts post-publication critique, \href{https://elifesciences.org/about/peer-review}{eLife} publishes reviewed preprints with public assessments, and \href{https://www.researchhub.com/}{ResearchHub} has gone as far as creating its own currency to pay reviewers. What is missing is the whole: verified expertise, real incentives to take part, and a visible, evolving record of where each paper stands, with methods, statistics, and conclusions assessed separately.

None of this is easy. Open commenting draws little real review: in one analysis of \href{https://www.biorxiv.org/}{bioRxiv} preprints, only 12 percent of readers' comments were full review reports\textsuperscript{\hyperlink{ref-malicki2021}{13}}. Crowds also herd: in a randomized experiment, a single early up-vote inflated final ratings by 25 percent on average\textsuperscript{\hyperlink{ref-muchnik2013}{14}}. A continuous system would have to verify who counts as an informed peer, reward unique and valuable input whether it challenges the crowd or confirms it, and hide the running score until each reviewer has given their own judgment. AI-generated reviews should be heavily moderated or banned; reviewers should write free-form critiques or redline the paper itself, which may help ensure they have read it in full. These are design problems, not reasons to keep a system whose label we already know we cannot fully trust and rely on.

A paper should not be peer reviewed once and stamped clean forever. It should stay under review for as long as it is relevant to the scientific enterprise, if not longer.

\vspace{14pt}
{\color{rule}\noindent\rule{0.3\textwidth}{0.4pt}}
\vspace{4pt}

{\fontsize{11.7}{16.38}\selectfont\setlength{\parindent}{0pt}\setlength{\parskip}{5pt}
{\itshape If you are working on this problem, or would like to discuss it further, please get in touch: \href{https://cortriatriatum.github.io/me/\#contact}{\upshape cortriatriatum.github.io/me}.}

Rick Mathews, PhD, is an MD/PhD candidate at the Oregon Health \& Science University (OHSU), currently on leave.

}

\vspace{12pt}
{\fontsize{10.5}{13.5}\selectfont\setlength{\parskip}{0pt}
\noindent{\bfseries References}\par\vspace{4pt}
\begin{enumerate}[leftmargin=1.8em,itemsep=5pt,topsep=0pt,label=\arabic*.]
\item \hypertarget{ref-fanelli2009}{}Fanelli D. How many scientists fabricate and falsify research? A systematic review and meta-analysis of survey data. \textit{PLoS ONE}. 2009;4(5):e5738. \url{https://doi.org/10.1371/journal.pone.0005738}\\[1pt]{\color{quote}\itshape ``A pooled weighted average of 1.97\% (N = 7, 95\%CI: 0.86-4.45) of scientists admitted to have fabricated, falsified or modified data or results at least once''}~{\color{quote}\footnotesize(Abstract)}
\item \hypertarget{ref-fang2012}{}Fang FC, Steen RG, Casadevall A. Misconduct accounts for the majority of retracted scientific publications. \textit{Proc Natl Acad Sci USA}. 2012;109(42):17028--17033. \url{https://doi.org/10.1073/pnas.1212247109}\\[1pt]{\color{quote}\itshape ``In contrast, 67.4\% of retractions were attributable to misconduct''}~{\color{quote}\footnotesize(Abstract)}
\item \hypertarget{ref-osc2015}{}Open Science Collaboration. Estimating the reproducibility of psychological science. \textit{Science}. 2015;349(6251):aac4716. \url{https://doi.org/10.1126/science.aac4716}\\[1pt]{\color{quote}\itshape ``Thirty-six percent of replications had statistically significant results''}~{\color{quote}\footnotesize(Abstract)}
\item \hypertarget{ref-schroter2008}{}Schroter S, Black N, Evans S, Godlee F, Osorio L, Smith R. What errors do peer reviewers detect, and does training improve their ability to detect them? \textit{J R Soc Med}. 2008;101(10):507--514. \url{https://doi.org/10.1258/jrsm.2008.080062}\\[1pt]{\color{quote}\itshape ``At baseline (Paper 1) reviewers found an average of 2.58 of the nine major errors''}~{\color{quote}\footnotesize(Abstract, Results)}
\item \hypertarget{ref-butler2023}{}Butler LA, Matthias L, Simard MA, Mongeon P, Haustein S. The oligopoly's shift to open access: how the big five academic publishers profit from article processing charges. \textit{Quant Sci Stud}. 2023;4(4):778--799. \url{https://doi.org/10.1162/qss_a_00272}\\[1pt]{\color{quote}\itshape ``we estimate that globally authors paid \$1.06 billion in publication fees to these publishers from 2015--2018.''}~{\color{quote}\footnotesize(Abstract)}
\item \hypertarget{ref-buranyi2017}{}Buranyi S. Is the staggeringly profitable business of scientific publishing bad for science? \textit{The Guardian}. June 27, 2017. \url{https://www.theguardian.com/science/2017/jun/27/profitable-business-scientific-publishing-bad-for-science}\\[1pt]{\color{quote}\footnotesize No verbatim excerpt attached.}
\item \hypertarget{ref-lariviere2015}{}Larivi\`{e}re V, Haustein S, Mongeon P. The oligopoly of academic publishers in the digital era. \textit{PLoS ONE}. 2015;10(6):e0127502. \url{https://doi.org/10.1371/journal.pone.0127502}\\[1pt]{\color{quote}\itshape ``Combined, the top five most prolific publishers account for more than 50\% of all papers published in 2013.''}~{\color{quote}\footnotesize(Abstract)}
\item \hypertarget{ref-aczel2021}{}Aczel B, Szaszi B, Holcombe AO. A billion-dollar donation: estimating the cost of researchers' time spent on peer review. \textit{Res Integr Peer Rev}. 2021;6:14. \url{https://doi.org/10.1186/s41073-021-00118-2}\\[1pt]{\color{quote}\itshape ``We found that the total time reviewers globally worked on peer reviews was over 100 million hours in 2020''}~{\color{quote}\footnotesize(Abstract, Results)}
\item \hypertarget{ref-koltun2021}{}Koltun V, Hafner D. The h-index is no longer an effective correlate of scientific reputation. \textit{PLoS ONE}. 2021;16(6):e0253397. \url{https://doi.org/10.1371/journal.pone.0253397}\\[1pt]{\color{quote}\itshape ``We find that the correlation of the h-index with awards that indicate recognition by the scientific community has substantially declined.''}~{\color{quote}\footnotesize(Abstract)}
\item \hypertarget{ref-ritter2017}{}Ritter S. Crowd-based peer review passes test. \textit{Chem Eng News}. 2017;95(24). \url{https://doi.org/10.1021/cen-09524-notw4}\\[1pt]{\color{quote}\itshape ``gave a stable of about 100 of the journal's reviewers 72 hours to respond anonymously to papers of their choosing''}~{\color{quote}\footnotesize(Article text)}
\item \hypertarget{ref-vangemmeren2020}{}van Gemmeren M, List B. How and why crowd reviewing works. \textit{Synlett}. 2020. \url{https://doi.org/10.1055/s-0040-1705955}\\[1pt]{\color{quote}\itshape ``it is now used as the default method for the evaluation of manuscripts at Synlett and SynOpen''}~{\color{quote}\footnotesize(Abstract)}
\item \hypertarget{ref-bik2016}{}Bik EM, Casadevall A, Fang FC. The prevalence of inappropriate image duplication in biomedical research publications. \textit{mBio}. 2016;7(3):e00809-16. \url{https://doi.org/10.1128/mBio.00809-16}\\[1pt]{\color{quote}\itshape ``Overall, 3.8\% of published papers contained problematic figures, with at least half exhibiting features suggestive of deliberate manipulation.''}~{\color{quote}\footnotesize(Abstract)}
\item \hypertarget{ref-malicki2021}{}Mali\v{c}ki M, Costello J, Alperin JP, Maggio LA. Analysis of single comments left for bioRxiv preprints till September 2019. \textit{Biochem Med}. 2021;31(2):020201. \url{https://doi.org/10.11613/BM.2021.020201}\\[1pt]{\color{quote}\itshape ``Twelve percent of non-author comments (N = 168) were full review reports traditionally found during journal review''}~{\color{quote}\footnotesize(Abstract, Results)}
\item \hypertarget{ref-muchnik2013}{}Muchnik L, Aral S, Taylor SJ. Social influence bias: a randomized experiment. \textit{Science}. 2013;341(6146):647--651. \url{https://doi.org/10.1126/science.1240466}\\[1pt]{\color{quote}\itshape ``created accumulating positive herding that increased final ratings by 25\% on average''}~{\color{quote}\footnotesize(Abstract)}
\end{enumerate}
}

\clearpage
{\fontsize{11.7}{16.38}\selectfont\setlength{\parindent}{0pt}\setlength{\parskip}{6pt}
\noindent{\bfseries How this was edited}\par
I wrote the original on November 17, 2026, and revised it the same day with editorial help from Claude, an AI model made by Anthropic. Claude critiqued the draft, found and checked the sources, and drafted edits, some on its own suggestion and some at my direction; I chose what to keep. The Markup view shows every change from the original. Every quote in the references is copied word for word from the published source, except where a reference says no excerpt is attached.\par
\vspace{10pt}
{\color{rule}\noindent\rule{\textwidth}{0.4pt}}\par
\noindent{\bfseries Original version: \textit{Peer Review Should be Scaled and Continuous for Every Paper}, November 17\textsuperscript{th}, 2026.} Unedited.\par
}
\fontsize{13.0}{18.85}\selectfont\setlength{\parskip}{6.5pt}
Formal peer-review as it is practiced today is significantly flawed and should be overhauled.

The magnitude of plagiarism, falsification, fabrication and lack of reproducibility is testament to the inability of the peer review system to independently verify the findings of scientific output as it is practiced today.

The guardian wrote a piece about how the publication system breaks the mechanisms of capitalism in sourcing unpaid labor and turning to the exact same consumer in sourcing paid demand for their own labor.

While services such as researchhub have gone so far as creating entire new currencies to reimagine this system, the bureaucracy of academia has bent to the exorbitant fees to access their own output as a field as well as retain professional progression to h indices, publication counts, and impact factors, none of which are true measures of true research impact.

For peer review at least, a scaled and continuous form seems to be an obvious solution applicable today. Scientific papers through preprint servers that may be able to invite comment and annotate consensus through a verifiably informed group of peers can make scientific papers into living, breathing documents with continuous peer review sourced from a scaled and an informed crowd vs the 2-3 peer reviewers editors struggle everyday to source who may not be the best suited to the task. The label it then provides for that paper is often unjustified.


\vspace{10pt}
{\centering\footnotesize \copyright\ 2026 Rick Mathews. All rights reserved.\par}
\end{document}
